\section{Phân tích yêu cầu hệ thống}
\label{sec:requirements}
\subsection{Ràng buộc và mức độ ưu tiên}
SRS v2.0 mô tả website doanh nghiệp kết hợp mua hàng cơ bản, với ba actor và bốn nhóm việc học phần. Các ràng buộc được dẫn lại gồm HTML/CSS/JavaScript, PHP phía server, MySQL, không framework PHP, TailAdmin Free, upload ảnh, validation hai phía, responsive, SEO, kiểm tra W3C và trình duyệt.\cite{srs}

Các FR dưới đây đều thuộc P0 của SRS; không đồng nghĩa tất cả đã được triển khai. Nguồn BT/TK được giữ để phân biệt yêu cầu học phần với chi tiết do dự án chọn. Khác biệt giữa SRS và mã nguồn được tập hợp tại phần đánh giá, đặc biệt PHP Views, TailAdmin, danh mục, mật khẩu và avatar.

\subsection{Actors và ma trận quyền}
\begin{table}[H]\centering
\caption{Ma trận quyền theo actor}\label{tab:permissions}
\begin{tabularx}{\linewidth}{@{}YP{2.0cm}P{2.5cm}P{2.5cm}@{}}\toprule
Chức năng & Guest & Member & Admin\\ \midrule
Xem trang công khai, tìm sản phẩm/tin & Có & Có & Có\\
Đăng ký, đăng nhập & Có & Đã có phiên & Đã có phiên\\
Quản lý giỏ & PHP Session & Giỏ MySQL & Giỏ cá nhân\\
Checkout COD và đơn cá nhân & Không & Đơn của mình & Đơn của mình\\
Hồ sơ, đổi mật khẩu, avatar & Không & Của mình & Của mình\\
Gửi bình luận & Không & Có & Có\\
Quản trị tài nguyên và trạng thái đơn & Không & Không & Có\\
\bottomrule\end{tabularx}
\end{table}

Ẩn menu trên giao diện không phải kiểm soát truy cập. Server phải kiểm role của session cho mọi thao tác quản trị và kiểm chủ sở hữu khi đọc đơn hàng. Mọi tham số như \code{admin=1}, \code{userId} hoặc ID trên URL phải được xem là đầu vào chưa tin cậy.

\subsection{Bảng yêu cầu chức năng}
{\small
\begin{longtable}{@{}P{1.8cm}P{9.6cm}P{1.3cm}P{1.8cm}@{}}
\caption{Yêu cầu chức năng P0 và liên kết kiểm thử}\label{tab:fr}\\
\toprule FR & Nội dung & Nguồn & TC\\\midrule\endfirsthead
\toprule FR & Nội dung & Nguồn & TC\\\midrule\endhead
FR-A01 & Đăng ký bằng họ tên, email, mật khẩu; không trùng email. & BT/TK & TC-01–02\\
FR-A02 & Đăng nhập email/mật khẩu, đăng xuất và kết thúc session. & BT/TK & TC-03\\
FR-A03 & Guest/Member/Admin; chặn sai quyền và tài khoản khóa. & BT/TK & TC-04–05\\
FR-A04 & Cập nhật họ tên, điện thoại, địa chỉ; đổi mật khẩu. & BT & TC-06\\
FR-A05 & Upload/thay avatar; file lưu trên server. & BT & TC-06\\
FR-A06 & Admin tìm kiếm/phân trang user; sửa, khóa/mở hoặc vô hiệu hóa. & BT & TC-07\\
FR-A07 & Admin reset mật khẩu an toàn, không xem mật khẩu gốc. & BT/TK & TC-07\\
FR-H01 & Trang chủ có nhận diện, menu, banner, danh mục, sản phẩm và tin mới. & BT/TK & TC-08\\
FR-H02 & Hiển thị và quản trị hotline, email, địa chỉ, logo, thông tin công ty. & BT & TC-08\\
FR-H03 & Guest/Member gửi liên hệ: họ tên, email, chủ đề, nội dung. & BT/TK & TC-09–10\\
FR-H04 & Admin tìm/phân trang liên hệ; đánh dấu unread/read/replied. & BT & TC-11\\
FR-H05 & Admin xóa liên hệ; sửa nội dung/banner trang chủ. & BT & TC-08,11\\
FR-F01 & Trang giới thiệu có nội dung, hình ảnh và thông tin hỗ trợ. & BT & TC-12\\
FR-F02 & Danh sách FAQ có thao tác mở/đóng trả lời. & BT/TK & TC-12\\
FR-F03 & Admin sửa nội dung và ảnh giới thiệu. & BT & TC-12\\
FR-F04 & Admin thêm/sửa/xóa/ẩn FAQ, sắp xếp hiển thị. & BT/TK & TC-12\\
FR-F05 & Hiển thị dịch vụ và bảng giá lắp ráp/vệ sinh/nâng cấp. & BT/TK & TC-12\\
FR-F06 & Admin thêm/sửa/ẩn/xóa dịch vụ và bảng giá. & BT/TK & TC-12\\
FR-P01 & Danh sách phân trang; chi tiết tên, ảnh, giá, mô tả, thông số. & BT/TK & TC-13\\
FR-P02 & Tìm sản phẩm theo từ khóa và lọc danh mục. & BT/TK & TC-13\\
FR-P03 & Admin CRUD/ẩn sản phẩm, sửa giá và tồn kho. & BT & TC-14–15\\
FR-P04 & Quản lý danh mục và upload ảnh sản phẩm. & BT/TK & TC-16\\
FR-P05 & Thêm giỏ, đổi số lượng, xóa sản phẩm và xem tổng tiền. & BT/TK & TC-17\\
FR-P06 & Giỏ Guest theo session, tiếp tục giỏ sau đăng nhập. & TK & TC-17\\
FR-P07 & Member nhập người nhận, đặt COD; server tính giá/kiểm kho. & TK & TC-18–19\\
FR-P08 & Lưu order/items MySQL, mã/trạng thái đơn, đơn cá nhân. & TK & TC-18,20\\
FR-P09 & Admin tìm/phân trang đơn, chuyển trạng thái hợp lệ. & BT & TC-21–22\\
FR-P10 & Admin tìm giỏ MySQL và đánh dấu trạng thái xử lý. & BT/TK & TC-31\\
FR-N01 & Danh sách tin có tìm kiếm và phân trang. & BT & TC-23\\
FR-N02 & Chi tiết tin có nội dung, ảnh và ngày đăng. & BT & TC-23\\
FR-N03 & Admin CRUD/ẩn tin; quản lý SEO title/description/keywords. & BT & TC-23\\
FR-N04 & Member bình luận; Guest phải đăng nhập. & BT/TK & TC-24\\
FR-N05 & Admin tìm/phân trang, duyệt/ẩn/xóa bình luận. & BT/TK & TC-24\\
FR-X01 & Dashboard quản trị sử dụng TailAdmin Free. & BT & Rà UI\\
FR-X02 & Mọi form validate bằng JavaScript và PHP. & BT & TC-10,15,26\\
FR-X03 & Danh sách quản trị có phân trang và phản hồi thao tác. & BT/TK & TC-30\\
FR-X04 & Ảnh logo/avatar/sản phẩm/bài viết/giới thiệu upload thật. & BT & TC-06,16\\
FR-X05 & Responsive điện thoại, tablet, desktop. & BT & TC-28\\
FR-X06 & Title/description, URL rõ nghĩa, alt và semantic HTML. & BT/TK & Rà SEO\\
\bottomrule\end{longtable}
}


\subsection{Bảng yêu cầu phi chức năng}
Các mốc 360/768/1280 px và mục tiêu phản hồi nhỏ hơn 3 giây là tiêu chí kỹ thuật được SRS lựa chọn để kiểm tra, không được trình bày là benchmark đã đo hoặc toàn bộ ngưỡng chấm điểm do giảng viên công bố.
{\small
\begin{longtable}{@{}P{1.8cm}P{5.1cm}P{6cm}P{1.6cm}@{}}
\caption{Yêu cầu phi chức năng và cách kiểm chứng}\label{tab:nfr}\\
\toprule NFR & Yêu cầu & Tiêu chí kiểm tra & TC\\\midrule\endfirsthead
\toprule NFR & Yêu cầu & Tiêu chí kiểm tra & TC\\\midrule\endhead
NFR-01 & Responsive tại 360, 768 và 1280 px. & Không tràn ngang; điều hướng/form sử dụng được. & TC-28\\
NFR-02 & Chạy trên Chrome, Firefox, Edge. & Smoke test P0, ghi phiên bản và lỗi từng trình duyệt. & TC-29\\
NFR-03 & Kiểm tra HTML/CSS bằng W3C. & Lưu kết quả; sửa lỗi nghiêm trọng trước nộp. & TC-29\\
NFR-04 & Validation ở client và PHP server. & Bỏ qua JS vẫn từ chối dữ liệu sai; không lưu dở dang. & TC-10,15\\
NFR-05 & Hạn chế SQL Injection bằng prepared statements. & Không nối input vào SQL; kiểm allowlist sort/filter. & TC-27\\
NFR-06 & Ngăn XSS, escape output và kiểm soát rich text. & Payload trong comment/contact không thực thi. & TC-25\\
NFR-07 & CSRF cho thao tác ghi, quyền kiểm tại server. & Thiếu token hoặc Member gọi admin phải bị từ chối. & TC-05,26\\
NFR-08 & Hash mật khẩu, session và upload an toàn. & Không lộ hash; chỉ loại ảnh/kích thước được phép. & TC-03,16\\
NFR-09 & SEO cơ bản trên các trang chính. & Kiểm HTML nhận về: title, meta, heading, alt, URL. & Rà SEO\\
NFR-10 & Phân trang và trạng thái danh sách rỗng. & Dữ liệu vượt một trang; đúng tổng và trang rỗng. & TC-13,30\\
NFR-11 & Mục tiêu phản hồi local nhỏ hơn 3 giây. & Đo trên cấu hình/dữ liệu xác định; chưa có số đo. & Đo local\\
NFR-12 & Cài đặt tái lập với SQL, seed, config mẫu, README. & Máy/môi trường sạch chạy được theo hướng dẫn. & Cài sạch\\
\bottomrule\end{longtable}
}


\subsection{Quy tắc nghiệp vụ}
\begin{longtable}{@{}P{1.5cm}P{13.6cm}@{}}
\caption{Quy tắc nghiệp vụ theo SRS}\label{tab:business-rules}\\
\toprule Mã & Quy tắc và tác động thiết kế\\\midrule\endfirsthead
\toprule Mã & Quy tắc và tác động thiết kế (tiếp)\\\midrule\endhead
BR-01 & Email duy nhất; lưu hash mật khẩu. UQ trong database xử lý cả trường hợp đăng ký đồng thời.\\
BR-02 & Tài khoản khóa không đăng nhập; PHP kiểm quyền Admin thay vì dựa vào route frontend.\\
BR-03 & Giá/tồn kho không âm; số lượng mua là số nguyên dương và không vượt khả dụng.\\
BR-04 & Không checkout giỏ rỗng; chỉ thành viên đã xác thực đặt được đơn.\\
BR-05 & Đọc lại giá, coupon, phí và tồn kho từ server tại thời điểm chốt đơn.\\
BR-06 & Đơn, chi tiết, trừ kho và xóa giỏ nằm trong cùng transaction; lỗi không được để dữ liệu dở dang.\\
BR-07 & pending → confirmed → shipping → completed; chỉ hủy từ pending/confirmed.\\
BR-08 & Khi hủy hợp lệ, hoàn kho đúng một lần; completed/cancelled là trạng thái cuối.\\
BR-09 & order\_items giữ snapshot tên, SKU và giá; sửa catalog không sửa lịch sử đơn.\\
BR-10 & Sản phẩm đã được tham chiếu trong đơn ưu tiên ẩn/ngừng bán, không xóa phá lịch sử.\\
BR-11 & Chỉ bài published và bình luận approved được hiển thị công khai.\\
BR-12 & Liên hệ có unread/read/replied; dịch vụ/bảng giá là dữ liệu CRUD, chưa có thanh toán dịch vụ.\\
\bottomrule\end{longtable}

\subsection{Điều kiện nghiệm thu}
Một chức năng được chấp nhận khi có luồng chính hoạt động, xử lý đầu vào sai và quyền truy cập, dữ liệu giữ sau tải lại, cùng chứng cứ kiểm thử trên cấu hình đã chốt. Hoàn thành CRUD phải bao gồm phản hồi thành công/thất bại và phân trang theo yêu cầu, không chỉ có form trên màn hình. Với upload phải tồn tại file hợp lệ ở server; với đơn hàng phải đối chiếu database, tổng tiền và thay đổi tồn kho.

Ngoài nghiệm thu chức năng, bản nộp cần schema/migration, dữ liệu công khai để demo, hướng dẫn chạy và báo cáo. Việc đủ số trang báo cáo không thay điều kiện hoàn thiện P0. Các kết quả W3C, responsive và nhiều trình duyệt chỉ được đánh dấu đạt sau khi chạy và lưu bằng chứng.
\clearpage
